\subsection*{Models and data}

\paragraph{Geneformer.} We used Geneformer V2-316M~\cite{theodoris2023transfer} (18 transformer
layers, 1,152 hidden dimensions, 18 attention heads per layer) from HuggingFace
(\texttt{ctheodoris/\allowbreak Geneformer}, subfolder \texttt{Geneformer-V2-316M}). Perturbation
data come from the Replogle genome-scale CRISPRi resource~\cite{replogle2022mapping} as a single
processed matrix spanning four cell lines (K562, RPE1, Jurkat, HepG2; 643,413 cells, 6,546 measured
genes, 2,023 perturbation targets plus non-targeting controls).

\paragraph{Control cohort.} The atlas was built from 2,000 non-targeting control cells sampled at
random (seed 42) from the 39,165 non-targeting cells in the resource. Because the control condition
is defined by perturbation label, this cohort spans all four lines: 591 Jurkat, 581 K562, 567 RPE1
and 261 HepG2 cells (mean 2,028 genes per cell, 4,056,351 total token positions). The atlas
dictionaries are therefore learned from a multi-cell-line control population, and we refer to them
throughout as the multi-line control atlas rather than attributing them to any single line. Cohort
composition for every analysis is recorded in \suptab{cohorts}. For the multi-tissue
comparison we additionally used 3,000 cells from Tabula Sapiens~\cite{tabula2022tabula} (1,000
immune cells across 43 cell types, 1,000 kidney cells, 1,000 lung cells), yielding 5,623,164 token
positions.

\paragraph{Cell-line assignment in downstream analyses.} Analyses that compare perturbed against
control cells draw both from the same line, so that a difference between the two cannot reflect a
difference in cell-line composition. Causal patching is likewise run within a single stated line.
Where an analysis uses the atlas dictionary on cells of one line, this is a deliberate transfer
setting and is labelled as such.

\paragraph{scGPT.} We used scGPT whole-human~\cite{cui2024scgpt} (12 transformer layers, 512 hidden
dimensions, 8 attention heads per layer), trained on approximately 33 million single-cell
transcriptomic profiles. scGPT ingests each cell as a pair of parallel token sequences:
gene-identity tokens drawn from a vocabulary of approximately 60,700 genes, and \emph{binned}
expression-value tokens (log1p-normalized counts discretized into 51 bins by default). Genes are
sorted by expression level (descending) and padded or truncated to a maximum sequence length of
1,200. Pre-training uses an attention-masking scheme (``gene prompt'' and ``value prompt'' modes)
that enables generative prediction of expression values conditional on gene identity, which we
refer to throughout as a generative or causal-attention objective to distinguish it from
Geneformer's bidirectional masked-token objective. Activations were extracted from the same 3,000
Tabula Sapiens cells, yielding 3,561,832 token positions per layer.

\paragraph{Input handling for each model.} The two models differ at the first layer, and this
difference propagates through all downstream analyses. For Geneformer the input is deterministic
given a cell: the per-gene expression vector is sorted, each gene receives a rank token, and no
continuous magnitude is passed to the model. For scGPT both the gene-identity token and its bin
index are embedded and summed before the transformer stack, so the model has direct access to a
coarse magnitude signal. Residual-stream activations therefore capture, at every layer, a
representation whose only expression signal is gene \emph{rank} for Geneformer, and one that
includes both rank-ordered gene identity and a binned magnitude for scGPT.

\subsection*{Activation extraction}

For both models we performed forward passes with PyTorch hooks registered at each transformer
layer's output (after the residual connection), collecting hidden states at every gene position.
Activations were stored as memory-mapped NumPy arrays (float32). Extraction used Apple Silicon MPS
acceleration. For Geneformer ($d{=}1{,}152$) this is 336.4~GB for 18 K562 layers (18.7~GB per
layer) and 103.6~GB for four Tabula Sapiens layers; for scGPT ($d{=}512$), approximately 82~GB for
12 layers.

For the perturbation analyses, where the quantity of interest is a per-cell summary rather than the
raw residual stream, we instead hooked the target layer directly, encoded the captured hidden state
through the layer's SAE on the accelerator, and reduced to per-cell statistics (mean activation,
firing rate, and mean square of each feature over that cell's gene positions) before transfer.
Cells were batched by sequence length. This keeps the cell as the retained unit of observation and
makes the per-cell feature matrices small enough to resample freely.

\subsection*{TopK sparse autoencoder architecture and training}

We used TopK SAEs~\cite{makhzani2013ksparse} parameterized by input dimension $d$. The encoder maps
$\mathbf{x} \in \mathbb{R}^{d}$ (centered by subtracting the training-set mean $\boldsymbol{\mu}$)
through a linear projection to a pre-activation vector in $\mathbb{R}^{4d}$, followed by TopK
sparsification retaining only the $k{=}32$ largest activations, to which a ReLU is applied. The
decoder reconstructs via $\hat{\mathbf{x}} = W_\text{dec}\mathbf{h}_\text{sparse} +
\mathbf{b}_\text{dec}$, where $W_\text{dec}$ columns are re-normalized to unit norm after each
gradient step. Training minimized reconstruction MSE; sparsity is enforced structurally by the TopK
operation rather than by an $L_1$ penalty.

For Geneformer ($d{=}1{,}152$): 4,608 features per layer, trained on a 1M-position subsample per
layer. For scGPT ($d{=}512$): 2,048 features per layer, trained on all 3,561,832 positions per
layer. Both used Adam, learning rate $3 \times 10^{-4}$, batch size 4,096, and 5 epochs. The atlas
SAEs were trained with random seed 42, which also fixes the training subsample; the sensitivity of
every reported quantity to this choice is characterized in the seed-stability analysis below.

\subsection*{Feature analysis and ontology annotation}

For each alive feature (activation frequency $> 0$ on 100K held-out positions) we identified the
top 20 genes by mean activation magnitude. Each feature was tested for enrichment against Gene
Ontology Biological Process~\cite{ashburner2000go}, KEGG pathways~\cite{kanehisa2000kegg}, Reactome
pathways~\cite{jassal2020reactome}, STRING protein--protein interactions~\cite{szklarczyk2023string},
and TRRUST TF--target relationships~\cite{han2018trrust}. TF--target priors from
DoRothEA~\cite{garcia-alonso2019dorothea} were used at the A+B+C confidence subset and at the union
of all confidence levels. Enrichment was assessed with a one-sided Fisher's exact test and
Benjamini--Hochberg FDR at $\alpha = 0.05$, with term sizes restricted to 5--500 genes.

\subsection*{Comparison against the principal subspace}

At each layer we computed the exact eigendecomposition of the activation covariance from a
500,000-position sample (40 evenly spaced contiguous blocks), giving all $d$ principal directions
and their eigenvalues. Three comparisons follow.

\emph{Reconstruction-matched capacity.} The truncated-SVD rank $k^{*}$ at which cumulative
eigenvalue mass equals the SAE's variance explained on held-out positions.

\emph{Cumulative projection.} For each unit-normalized decoder direction $\mathbf{d}_f$ we computed
$\rho_k = \lVert P_k \mathbf{d}_f \rVert^2 = \sum_{i \le k} (\mathbf{d}_f \cdot \mathbf{v}_i)^2$,
the fraction of the direction's norm lying in the rank-$k$ principal subspace, for all $k$ from 1
to $d$. The reference for these curves is a random unit direction, for which
$\mathbb{E}[\rho_k] = k/d$.

\emph{Single-axis alignment.} The fraction of features whose absolute cosine similarity with at
least one of the leading $k$ axes exceeds $\tau$, reported jointly over
$k \in \{1, \dots, d\}$ and $\tau \in \{0.3, 0.5, 0.7, 0.9\}$ so that the figure does not depend on
either choice.

\emph{Annotation yield per direction.} The leading 50 principal axes were annotated with the
identical top-20-gene enrichment protocol used for SAE features. Because a principal axis is signed
whereas a TopK feature is one-sided, each axis was annotated in both polarities (genes with the
most positive and most negative mean projection), and results are reported per polarity and for
their union.

\subsection*{Counting distinct concepts across layers}

Dictionary atoms from different layers occupy different vector spaces and cannot be pooled as
directions. They can be compared in the shared gene space through their top-20 gene signatures. We
built a similarity graph over all 82,525 Geneformer features, placing an edge between two features
whose top-20 gene sets share at least 5 genes and have Jaccard similarity at least 0.14, and
partitioned it with the Leiden algorithm~\cite{traag2019louvain} at resolutions
$\gamma \in \{0.5, 1.0, 2.0\}$. The number of communities is the number of distinct gene-level
programs. We additionally report, per layer, the number of atoms, the number of distinct programs,
the mean absolute decoder coherence, and the Welch bound
$\sqrt{(n-d)/(d(n-1))}$ for $n$ unit vectors in $\mathbb{R}^{d}$, which is the smallest achievable
maximum coherence and therefore the correct reference for how many near-orthogonal directions a
$d$-dimensional space can host.

\subsection*{Co-activation graph and module detection}

For each layer, pointwise mutual information (PMI~\cite{manning1999foundations}) was computed
between all pairs of alive features across the layer's token positions; feature $i$ is active at
position $j$ if it is among the top-$k$ activations there, and
$\text{PMI}(i,j) = \log_2 P(i,j) / [P(i)P(j)]$. Edges were retained above a significance threshold
($p < 0.001$ under a 1,000-sample permutation null in which feature activations were independently
column-shuffled). Community detection used the Leiden algorithm~\cite{traag2019louvain} at
resolution $\gamma = 1.0$, the \texttt{leidenalg} default. Because a single resolution is a choice
rather than a ground truth, we swept $\gamma \in \{0.5, 0.75, 1.0, 1.5, 2.0\}$ at layers 0 and 11
and report the resulting partition sizes and adjusted Rand indices; module counts in the main text
are the $\gamma{=}1.0$ cut of a hierarchical structure whose qualitative module identities are
recognizable throughout the range.

\subsection*{Causal feature patching}

Single-feature ablation was performed with PyTorch forward hooks. At the target layer's output the
hidden state was encoded through the SAE, one feature's activation was zeroed, the state was
decoded and substituted, and the forward pass continued to the output logits. Disruption is the
mean change in output logits at a set of gene positions.

Three properties of the design keep the evaluation independent of the annotation. First, features
are drawn from three arms: the most richly annotated features, a uniformly random sample of
annotated features, and a uniformly random sample of all alive features. Second, for each feature we
partition gene positions into those whose gene is in both the feature's top-20 activating genes and
its annotated term, those in the annotated term but \emph{not} in the top-20 (the independent
evaluation set), and all others. Third, for every feature we draw a random gene set matched in size
and in mean-activation decile to the independent evaluation set, giving each feature its own null.
Specificity is the ratio of mean absolute disruption on a gene set to that on all other positions,
and the held-out ratio is reported against the matched-random ratio for the same feature.

\subsection*{Perturbation response mapping}

For each perturbation target we sampled up to 200 CRISPRi cells and compared them against 600
non-targeting control cells of the same line, encoding every cell through the layer-11 SAE as
described above. A feature responds if its per-cell mean activation differs between perturbed and
control cells (Mann--Whitney $U$ over cells, Benjamini--Hochberg across features, FDR $< 0.05$) with
$|d| > 0.5$ (Cohen's $d$ over cells). The cell is the unit of replication throughout: gene positions
within a cell are not independent observations, and we report the intraclass correlation of feature
activation within cells together with the number of features that a position-pooled test would call
significant on exactly the same cells.

TF specificity is evaluated in the same framework as every other method (below): the union of the
top-20 gene lists of responding features, ranked by effect size, forms the predicted target set,
which is tested for enrichment of the TF's curated targets by hypergeometric test with
Benjamini--Hochberg correction across the panel. A label-permutation null, in which perturbed and
control labels are shuffled across the same cells, calibrates the procedure.

\subsection*{Recoverability benchmark}

A recovery rate is interpretable only against the amount of information the measurements contain.
We therefore evaluated several methods on one footing: each emits a ranked list of putative targets
for every TF in the panel, the top 100 of which are tested for enrichment of that TF's curated
targets by hypergeometric test with Benjamini--Hochberg correction across the panel.

\emph{Panel.} TFs perturbed in the cell line with at least 50 cells and at least 5 curated targets
present among the 6,546 measured genes. The second criterion is reported explicitly, because a TF
whose curated targets are not measured cannot show target-specific behaviour under any method.

\emph{Universe.} The hypergeometric universe is the set of measured genes. Using a nominal
20,000-gene universe is anticonservative here, since every method can only draw predictions from
measured genes; we report both, and the random-gene-set control is calibrated only under the
measured-gene universe.

\emph{Methods.} (i)~Differential expression of the knockdown itself (Mann--Whitney over cells,
perturbed vs non-targeting, Benjamini--Hochberg across genes), which is perturbation-aware and
constitutes the empirical ceiling attainable on these data. (ii)~GENIE3~\cite{huynhthu2010genie3}
and (iii)~GRNBoost2~\cite{moerman2019grnboost2}, in their published formulation: every target gene
is regressed on the candidate regulators with a tree ensemble and regulators are ranked by variable
importance (random forest with \texttt{max\_features}~=~$\sqrt{p}$; and stochastic gradient boosting
with learning rate 0.01, \texttt{max\_features}~=~0.1, subsample 0.9 and early stopping), fitted on
the same control cells. (iv)~Pearson and (v)~Spearman co-expression with the TF on the same control
cells. (vi)~Cosine similarity between mean contextual gene embeddings taken from the Geneformer
residual stream, a model-internal baseline that uses no SAE. (vii)~Size-matched random gene sets.
(viii)~The SAE feature response described above.

\subsection*{Replication across cell lines}

The perturbation analysis was repeated in RPE1, Jurkat and HepG2 cells from the same CRISPRi
resource, which are independent screens in independent cell backgrounds processed identically. In
every line, perturbed cells are compared against non-targeting cells of that same line. Each line
was analysed with two dictionaries: the multi-line control atlas dictionary, and the dictionary
additionally trained on Tabula Sapiens primary tissue, which asks whether widening the training
distribution beyond immortalized lines changes what is recovered.

\subsection*{Run-to-run stability}

The atlas SAE was retrained at layers 0, 5, 11 and 17 under five independent seeds with the
training subsample held fixed, isolating initialization and batch-order stochasticity, plus one
additional run per layer with a different 1M-position training subsample. Dictionaries were compared
by decoder cosine similarity, reporting both the unconstrained best match for each atom and a greedy
one-to-one assignment. Variance explained, dead-feature count, decoder coherence, annotation rate,
and module count were recomputed per run.

\subsection*{Matched cross-model comparison}

Because the two atlases were built from different inputs, architecture and input distribution are
confounded in any direct comparison of their statistics. We therefore trained Geneformer SAEs on the
same Tabula Sapiens cells used for the scGPT atlas, with identical hyperparameters, and compared the
two models on matched data at matched relative depth (layer index divided by depth minus one). For
the Geneformer arm we additionally separate two effects by reporting the variance explained when the
K562-trained SAE is applied to Tabula Sapiens activations against that of the SAE trained on those
activations directly.

\subsection*{Interactive web atlases}

Both feature atlases were deployed as single-page web applications built with React 18, TypeScript,
Vite 6, Tailwind CSS, and Plotly.js. All feature data were preprocessed into compact JSON files
served statically via GitHub Pages.

\subsection*{Dimensionality reduction and visualization}

We projected decoder weight vectors to two dimensions using UMAP~\cite{mcinnes2018umap} with cosine
distance, obtaining structureless point clouds as expected from quasi-orthogonal directions. We
therefore visualized the co-activation graph directly with a force-directed layout
(Fruchterman--Reingold via NetworkX); the corresponding supplementary figures are force-directed
layouts of that graph and are not UMAP or t-SNE projections. For independent validation,
annotation-based feature vectors (TF-IDF weighted ontology terms) were projected with UMAP
($n_\text{neighbors}=15$, $\min_\text{dist}=0.1$, cosine metric) and t-SNE (perplexity 20, cosine
metric).

\subsection*{Computational environment}

All experiments were run on a single Apple Silicon workstation (Apple M2 Pro, 10 cores, 32~GB
unified memory). Acceleration used PyTorch's MPS backend, with CPU fallback for operations without
an MPS implementation (Leiden community detection, ontology Fisher tests, tree-ensemble regression).
Software stack: Python 3.11; PyTorch 2.x with MPS; NumPy, SciPy, scikit-learn, pandas, scanpy,
anndata, leidenalg, python-igraph, networkx, umap-learn, statsmodels, h5py. Exact package versions
are pinned in \texttt{requirements.txt} in the code repository. Wall-clock times for the longest
steps are logged alongside the compute environment: Geneformer K562 activation extraction
$\sim$8~h, SAE training $\sim$3~min per layer, per-layer PMI and Leiden $\sim$20--60~min,
per-cell perturbation encoding $\sim$0.9 cells~s$^{-1}$.
